%! Polynomial Functions and Rates of Change — Unit 2, Lesson 2.2 — Lesson Plan
\documentclass[10pt]{article}
\usepackage{precalculus-article}
\usepackage{precalculus-boxes}
\newcommand{\TallMath}[1]{$\displaystyle #1\rule[-1.4em]{0pt}{3.2em}$}
\newcommand{\CourseName}{Precalculus}
\newcommand{\MeetingLength}{55 minutes}
\newcommand{\UnitNumberName}{Unit 2: Polynomial Functions \quad}
\newcommand{\LessonNumberName}{Lesson 2.2: Polynomial Functions and Rates of Change}

\begin{document}
\begin{center}
    {\Large\bfseries \CourseName \\
    {\small\bfseries \UnitNumberName \LessonNumberName}}
\end{center}

\begin{tcolorbox}[colback=lilac, colframe=plum, boxrule=0.9pt, arc=2mm,
    left=2mm, right=2mm, top=1.5mm, bottom=1.5mm]
    \textbf{Primary Objective:} Students will read the increasing and
    decreasing intervals of a polynomial function from a graph and from a
    table; identify local (relative) and global (absolute) maxima and minima,
    including at an included endpoint; justify why a turn must occur between
    two distinct real zeros and why even degree forces a global extremum; and
    locate a point of inflection as the input where the \textit{rate} of
    change turns around. \hfill\textit{\small (CED Topic 1.4 --- Skills 2.A, 3.A)}
\end{tcolorbox}

\renewcommand{\arraystretch}{1.6}
\begin{skillbox}[Priority Ideas \& Skills]{goldbox}
\begin{tabularx}{\linewidth}{@{}X X@{}}
\textbf{AP Skills} &
\textbf{Key Understandings} \\
\midrule
\textbf{2.A} --- Identify information from graphical, numerical, analytical,
or verbal representations to answer a question or construct a model. &
\begin{itemize}[leftmargin=1.2em,itemsep=2pt,topsep=0pt]
  \item A nonconstant polynomial is any expression equivalent to
        $a_n x^n + \cdots + a_0$ with $n$ a positive integer and
        $a_n \ne 0$; a constant has degree zero. (1.4.A.1)
  \item Where a polynomial switches between increasing and decreasing --- or
        at an included endpoint --- it has a local maximum or minimum. The
        greatest local max is the global max; the least local min is the
        global min. (1.4.A.2)
\end{itemize} \\
\textbf{3.A} --- Describe the characteristics of a function with varying
levels of precision, depending on the representation available. &
\begin{itemize}[leftmargin=1.2em,itemsep=2pt,topsep=0pt]
  \item Between every two distinct real zeros of a nonconstant polynomial
        there is at least one local maximum or minimum. (1.4.A.3)
  \item Even-degree polynomials have a global maximum or a global minimum;
        odd-degree ones have neither. (1.4.A.4)
  \item A point of inflection is an input where the rate of change switches
        between increasing and decreasing --- equivalently, where concavity
        flips. (1.4.A.5)
\end{itemize} \\
\end{tabularx}
\end{skillbox}

\begin{skillbox}[{Vocabulary, Concepts, \& Theorems}]{greenbox}
\renewcommand{\arraystretch}{1.4}
\begin{tabularx}{\linewidth}{@{} >{\leavevmode\bfseries\color{plum}}p{0.27\linewidth} X @{}}
Increasing / decreasing interval & An interval on which outputs rise as inputs rise, or fall as inputs rise. \\
Local (relative) maximum / minimum & An output value where the function switches between increasing and decreasing, or the value at an included endpoint of a restricted domain. \\
Global (absolute) maximum / minimum & The greatest of all local maxima; the least of all local minima --- the extreme output over the whole domain. \\
Concavity (from 2.1) & AROCs increasing $\Rightarrow$ concave up; AROCs decreasing $\Rightarrow$ concave down. \\
Point of inflection & An input where the rate of change switches between increasing and decreasing --- where concavity changes. \\
Zero of a polynomial & An input where the output is 0; today, only a landmark that forces a turn between consecutive zeros. \\
\end{tabularx}
\end{skillbox}

\begin{skillbox}[Activate Prior Knowledge \& Spiral Review]{lilac}
    Please see attached warm-up (spiral review).
    Skills reviewed: naming degree, leading term, and leading coefficient from
    standard form (Lesson 2.0); computing an average rate of change over an
    interval from a table (Lessons 1.3 and 2.1); and reading concavity off a
    list of AROCs over consecutive equal-length intervals (Lesson 2.1). Item
    3's ``decreasing AROCs $\Rightarrow$ concave down'' is the exact move
    Part 4 rebuilds into the definition of a point of inflection.
\end{skillbox}

\begin{skillbox}[Hook]{lilac}
    \textbf{Two summits, one highest point.} A trail guide lists two summits
    on the Ridge Trail: Bell Knob at 1{,}200 feet and Signal Peak at 1{,}850
    feet.
    \begin{enumerate}[itemsep=2pt]
        \item Ask: are both of these really summits? (Yes --- each is the
              high point of its own stretch of trail.)
        \item Then: which one would you put on the sign at the trailhead?
              (Signal Peak --- it is the highest point \textit{anywhere} on
              the trail.)
    \end{enumerate}
    \textit{Discussion prompt:} mathematics gives those two ideas different
    names. Bell Knob is a \textbf{local} maximum; Signal Peak is a local
    maximum \textit{and} the \textbf{global} maximum. Today every graph gets
    read that way.
\end{skillbox}

\begin{skillbox}[Lesson]{lilac}
\begin{multicols}{2}

\textbf{Part 1: The Definition, Stated Precisely} (7 min)

\begin{itemize}
  \item Restate 1.4.A.1 once, cleanly: $p(x) = a_n x^n + \cdots + a_0$, $n$ a
        positive integer, coefficients real, $a_n \ne 0$. Degree, leading
        term, leading coefficient. A constant is degree zero.
  \item This is review from 2.0 --- keep it to the quick check
        ($f(x) = 4x^3 - x^2 + 9$) and move on. The lesson is about
        \textit{behavior}, not form.
\end{itemize}

\textbf{Part 2: Increasing, Decreasing, Local Extrema} (13 min)

\begin{itemize}
  \item Project the notes' Section 2 graph on $[0,6]$ with points $A$--$D$
        pre-marked. Students only read and label --- nothing is drawn.
  \item Increasing on $(0,2)$ and $(4,6)$; decreasing on $(2,4)$. $B$ is a
        local max (value 6 at $x=2$); $C$ a local min (value 2 at $x=4$).
  \item Then the endpoint clause of 1.4.A.2: the domain is closed, so $A$ is
        a local minimum and $D$ a local maximum \textit{because they are
        included endpoints}. Students almost never volunteer this.
  \item Language discipline: the max is the \textit{value}; it
        \textit{occurs at} an input.
\end{itemize}

\columnbreak

\textbf{Part 3: Local vs.\ Global, and What Zeros Force} (15 min)

\begin{itemize}
  \item Same graph: list local maxima $\{6, 7\}$ and local minima
        $\{1, 2\}$. Greatest and least give the global max 7 (at $x=6$) and
        global min 1 (at $x=0$). $B$ is the trap --- a local max that is not
        the global one.
  \item Section 4's graph has zeros at $1, 3, 5$. Ask what the curve
        \textit{had} to do between consecutive zeros; land on 1.4.A.3.
        Generalize: $k$ distinct real zeros force at least $k-1$ turns.
  \item Even degree $\Rightarrow$ a global max or a global min (1.4.A.4);
        odd degree $\Rightarrow$ neither, because the two ends run opposite
        directions. Stay verbal --- formal end behavior is Lesson 2.5.
\end{itemize}

\textbf{Part 4: Points of Inflection} (15 min)

\begin{itemize}
  \item Return to 2.1's tool: AROCs over consecutive equal-length intervals.
        The notes' table gives $2, 5, 9, 8, 4$ --- rising, then falling.
  \item The \textit{rate} turns around near $x=3$: concave up before,
        concave down after. That input is the point of inflection.
  \item Drive the contrast: every AROC in that table is positive, so the
        function never turns around at all --- no local extrema, yet an
        inflection point. Extremum = AROCs change \textit{sign};
        inflection = AROCs change \textit{direction}.
\end{itemize}
\end{multicols}
\end{skillbox}

\begin{skillbox}[Active Monitoring]{redbox}
While students work on guided notes and practice, circulate and check:
\begin{itemize}
  \item \textbf{Common error:} naming an extremum by its input instead of its
        output (``the maximum is 2''). Ask: ``Is 2 the height or the
        location?''
  \item \textbf{Common error:} calling every local maximum the global
        maximum. Point at $B$: ``Is anything on this graph higher?''
  \item \textbf{Common error:} confusing an inflection point with a turning
        point. Diagnostic question: ``Did the function turn around there, or
        just the rate?''
  \item \textbf{Common error:} placing a required extremum \textit{at} a
        zero rather than \textit{between} two of them.
  \item \textbf{Watch for:} students skipping the included endpoints when
        listing local extrema on a restricted domain.
\end{itemize}
\end{skillbox}

\begin{skillbox}[Group Work \& Differentiation]{redbox}
Students work on the \textbf{Group Activity: The Drone Survey} in leveled
tiers. Every tier uses the same table and pre-drawn altitude graph on
$[0,6]$ --- no tier draws anything.

\begin{itemize}
  \item \textbf{Tier R (Remediate):} Read climbing/descending intervals and
        the four labeled points off the graph; classify $Q$ and $R$ as local
        extrema and $P$ and $S$ as endpoint extrema; report the global max
        (500 ft) and global min (100 ft).
  \item \textbf{Tier A (Approaching Proficiency):} From the table alone,
        fill the AROC row ($200, -160, -140, 110, 170, 120$), locate the
        local max at $x=1$ and local min at $x=3$ from sign changes, then
        find the two inflection points from where the AROCs bottom out
        ($x \approx 2$) and top out ($x \approx 5$).
  \item \textbf{Tier E (Extension):} At most $n-1$ local extrema for degree
        $n$; four distinct zeros force at least three turns, so degree
        $\ge 4$; justify 1.4.A.4 for even and odd degree; then the
        true/false item --- exactly one local extremum does \textit{not}
        force degree 2 ($x^4$ is the counterexample).
\end{itemize}
\end{skillbox}

\begin{skillbox}[Individual Work \& Assessment]{redbox}
\textbf{Exit Ticket} (last 5--7 minutes, individual, collected):
\begin{enumerate}
  \item From a pre-drawn graph on $[0,5]$: intervals of increase and
        decrease, and the inputs of the local maximum and local minimum.
  \item From the same graph: the global maximum and global minimum values and
        where they occur.
  \item From an AROC list ($3, 7, 12, 9, 5$): decide there is no local
        extremum, and locate the point of inflection and the concavity
        change.
\end{enumerate}
\textbf{Assessment note:} Item 2 separates local from global --- a student
who names 4 (the local max of item 1) as the global maximum has the exact
misconception the hook targeted. Item 3 needs no graph at all; a miss there
is a miss on the lesson's hardest idea, not on arithmetic.
\end{skillbox}

\begin{skillbox}[Reinforcement \& Extension]{goldbox}
\textbf{Homework:} 5 problems --- reading all four kinds of extrema off a
pre-drawn graph with included endpoints, an AROC row from a table with an
inflection point but no extremum, a minimum-extrema count from four distinct
zeros, three must-be-true statements on degree and zeros, and one AP-style
multiple-choice item on what an AROC list guarantees.

\textbf{Extension (optional):} Degree 5 with five distinct real zeros ---
the ``at least $k-1$'' and ``at most $n-1$'' bounds both equal 4, so the
count is pinned to exactly 4 local extrema.

\textbf{Preview:} Next lesson (2.3, Polynomial Functions and Real Zeros)
promotes the zeros from landmarks to the main event: finding them from a
factored form and describing what the graph does at each one.
\end{skillbox}

\begin{teachernote}[Warm-Up]
Five minutes, silent start, no calculators. Item 1 is pure 2.0 recall and
should take seconds; if it does not, the whole class needs a two-minute
standard-form reset before Part 1. Item 2's AROC is the computation Tier A of
the activity runs six times, so watch the subtraction order here. Item 3 is
the load-bearing one: ``decreasing AROCs $\Rightarrow$ concave down'' is
re-used verbatim in Part 4 to define a point of inflection --- note anyone who
guesses.
\end{teachernote}

\begin{teachernote}[Guided Notes]
Part 1 is deliberately thin --- resist re-teaching standard form. The whole
lesson lives in Sections 2--5. The included-endpoint clause in Section 2 is
the piece most classes skip and most AP items reward; say out loud that $A$
and $D$ are local extrema \textit{only because the domain stops there}.
Section 5 is the hard one: work the AROC row $2, 5, 9, 8, 4$ on the board
before students fill it in, and keep returning to the sentence ``every one of
these is positive, so the function never turned around.'' If time runs short,
compress Section 3 into Section 2's discussion --- do not cut Section 5, it is
exit ticket item 3 and homework items 2 and 5.
\end{teachernote}

\begin{teachernote}[Group Activity]
All three tiers share the drone table and graph, so groups can be re-sorted
mid-period without new setup. Tier A should keep the graph covered through
question 4 --- the point is that a table alone locates extrema and inflection
points; question 5 is the reveal. Tier A's two inflection answers
($x \approx 2$ and $x \approx 5$) are approximate by design: accept any
answer that names the right interval and says why. Tier E question 4 is the
best whole-class share in the closing minutes --- ``at most $n-1$'' is a
ceiling, not a count, and $x^4$ settles it in one line.
\end{teachernote}

\begin{teachernote}[Exit Ticket]
Collected, completion credit only. Scan item 2 first: a student who writes 4
as the global maximum is reading ``the tallest bump'' instead of ``the
greatest output,'' which is the endpoint idea from Section 2 --- sort those
for a two-minute reteach. On item 3(a) accept any reason equivalent to ``all
the AROCs are positive.'' On 3(b) accept $x = 3$ or ``between 2 and 4.''
\end{teachernote}

\begin{teachernote}[Homework]
Problem 1(c) is the included-endpoint check and problem 2 is the
inflection-without-extremum check; those two carry the diagnostic weight.
Problem 5 is the AP-style item: distractor A is the student who reads a
falling AROC as a falling function, and distractor D is the same error stated
directly --- both are worth two minutes of autopsy tomorrow. The extension is
the only place the two counting bounds are squeezed together; if Tier E never
got there, assign it explicitly and open 2.3's warm-up window with it.
\end{teachernote}
\end{document}
